\documentclass[tikz,border=20pt]{standalone}
\usetikzlibrary{calc, arrows.meta, shadings, shapes.geometric}
\usepackage[UTF8]{ctex}

\begin{document}
\begin{tikzpicture}[
    y=0.7cm, x=2.5cm,
    crit node/.style={circle, draw=black, fill=gray!20, thick, inner sep=1pt, minimum size=16pt, font=\sffamily\footnotesize},
    reg node/.style={circle, draw=black, fill=white, inner sep=0.5pt, minimum size=12pt, font=\sffamily\tiny},
    main edge/.style={very thick, draw=black!80},
    % ===== 中部：映射 =====
    axis/.style={line width=1.6pt, draw=black!85},
    proj/.style={draw=black!35, dashed, line width=0.7pt, -{Stealth[length=2mm, width=1.6mm]}},
    mapdot/.style={ellipse, fill=black!75, text=white, line width=1pt, inner sep=1.2pt,
                   minimum width=19pt, minimum height=13pt, font=\sffamily\scriptsize\bfseries},
    ilab/.style={font=\sffamily\footnotesize, text=black!80},
    idx/.style={font=\sffamily\tiny\bfseries, anchor=south, fill=white, inner sep=0.6pt},
    echi/.style={font=\sffamily\tiny, text=black!65, fill=white, inner sep=0.5pt},
    tagt/.style={font=\sffamily\tiny, text=red!70!black},
    tags/.style={font=\sffamily\tiny, text=green!45!black},
    % ===== 右部：等值面 =====
    iso3d/.style={very thick, draw=green!45!black},
    lab/.style={font=\sffamily\small, text=black!75},
]

% =========================================================
% ===== 左部：增广轮廓树（保持不变）=====
% =========================================================
\node[crit node] (C24) at (0, 24) {24};
\node[crit node] (C22) at (1, 22) {22};
\node[crit node] (C15) at (0.5, 15) {15};
\node[crit node] (C13) at (2, 13) {13};
\node[crit node] (C10) at (1, 10) {10};
\node[crit node] (C8)  at (2, 8) {8};
\node[crit node] (C1)  at (0, 1) {1};
\node[crit node] (C0)  at (0.5, 0) {0};

\draw[main edge] (C15) -- (C24);
\foreach \v in {17, 18, 20, 23} \path (C15) -- (C24) node[pos={(\v-15)/(24-15)}, reg node] {\v};

\draw[main edge] (C15) -- (C22);
\foreach \v in {16, 19, 21} \path (C15) -- (C22) node[pos={(\v-15)/(22-15)}, reg node] {\v};

\draw[main edge] (C1) -- (C15);
\foreach \v in {3, 6, 7, 14} \path (C1) -- (C15) node[pos={(\v-1)/(15-1)}, reg node] {\v};

\draw[main edge] (C10) -- (C15);
\foreach \v in {11, 12} \path (C10) -- (C15) node[pos={(\v-10)/(15-10)}, reg node] {\v};

\draw[main edge] (C10) -- (C13);
\draw[main edge] (C8) -- (C10);

\draw[main edge] (C0) -- (C10);
\foreach \v in {2, 4, 5, 9} \path (C0) -- (C10) node[pos={(\v-0)/(10-0)}, reg node] {\v};

% ---- 每条边(连通分量)的 Euler 示性数 χ 标注（初值 +1，鞍点后仍为 +1）----
\path (C0)  -- (C10) node[pos=0.68, echi] {$\chi{=}1$};
\path (C1)  -- (C15) node[pos=0.88, echi] {$\chi{=}1$};
\path (C8)  -- (C10) node[pos=0.5,  echi] {$\chi{=}1$};
\path (C10) -- (C13) node[pos=0.5,  echi] {$\chi{=}1$};
\path (C10) -- (C15) node[pos=0.5,  echi] {$\chi{=}1$};
\path (C15) -- (C22) node[pos=0.62, echi] {$\chi{=}1$};
\path (C15) -- (C24) node[pos=0.62, echi] {$\chi{=}1$};

% 左部标题
\node[font=\sffamily\bfseries\small, text=black!70, align=center] at (1.0, 25.8)
   {轮廓树\\{\footnotesize\mdseries（临界点 + 普通节点）}};

% =========================================================
% ===== 中部：映射（临界点 -> 标量轴分区）=====
% 每个极值点 / 鞍点沿等高度水平投影到标量轴上的一个点，
% 这些点把标量轴分成若干区间，每个区间内连通分量数 c 恒定。
% =========================================================

% ---- 区间色带（背景）：按连通分量数 c 着色，c 相同用同一颜色 ----
% c=1 -> blue!12 ；c=2 -> green!14 ；c=3 -> orange!20
% (0,1):c1  (1,8):c2  (8,10):c3  (10,13):c3  (13,15):c2  (15,22):c2  (22,24):c1
\fill[blue!12]   (3.78,0)  rectangle (4.78,1);   % I1 c=1
\fill[green!14]  (3.78,1)  rectangle (4.78,8);   % I2 c=2
\fill[orange!20] (3.78,8)  rectangle (4.78,10);  % I3 c=3
\fill[orange!20] (3.78,10) rectangle (4.78,13);  % I4 c=3
\fill[green!14]  (3.78,13) rectangle (4.78,15);  % I5 c=2
\fill[green!14]  (3.78,15) rectangle (4.78,22);  % I6 c=2
\fill[blue!12]   (3.78,22) rectangle (4.78,24);  % I7 c=1
\draw[black!30, line width=0.5pt] (3.78,0) rectangle (4.78,24);
\foreach \v in {1,8,10,13,15,22}
   \draw[black!30, line width=0.4pt] (3.78,\v) -- (4.78,\v);

% ---- 区间标注：区间编号 + 亏格 g（以初始亏格 g1 为基准）----
% 亏格 g = c - x/2，x 为每条边的 χ = 1（常数）。设 I1 亏格为 g1，
% 则各区间 g = g1 + (c-1)：I1,I7 -> g1；I2,I5,I6 -> g1+1；I3,I4 -> g1+2。
\node[ilab, align=center] at (4.28,0.5)  {$I_1$\\[-1pt]{\scriptsize$g{=}g_1$}};
\node[ilab, align=center] at (4.28,4.5)  {$I_2$\\[-1pt]{\scriptsize$g{=}g_1{+}1$}};
\node[ilab, align=center] at (4.28,9)    {$I_3$\\[-1pt]{\scriptsize$g{=}g_1{+}2$}};
\node[ilab, align=center] at (4.28,11.5) {$I_4$\\[-1pt]{\scriptsize$g{=}g_1{+}2$}};
\node[ilab, align=center] at (4.28,14)   {$I_5$\\[-1pt]{\scriptsize$g{=}g_1{+}1$}};
\node[ilab, align=center] at (4.28,18.5) {$I_6$\\[-1pt]{\scriptsize$g{=}g_1{+}1$}};
\node[ilab, align=center] at (4.28,23)   {$I_7$\\[-1pt]{\scriptsize$g{=}g_1$}};

% ---- 标量轴（线段，与树同高度）----
\draw[axis, -{Stealth[length=3mm, width=2.4mm]}] (3.6,-0.4) -- (3.6,25.2);
\node[font=\sffamily\small\itshape, text=black!70] at (3.6,25.7) {$f$};

% ---- 投影：每个临界点 -> 轴上一点（同高度水平映射）----
\foreach \n in {C0,C1,C8,C10,C13,C15,C22,C24} {}
\draw[proj] (C24) -- (3.47,24);
\draw[proj] (C22) -- (3.47,22);
\draw[proj] (C15) -- (3.47,15);
\draw[proj] (C13) -- (3.47,13);
\draw[proj] (C10) -- (3.47,10);
\draw[proj] (C8)  -- (3.47,8);
\draw[proj] (C1)  -- (3.47,1);
\draw[proj] (C0)  -- (3.47,0);

% ---- 轴上的临界点标记（深色椭圆；边框按类型着色）----
% 极小值 i=0（出生）：0,1,8 ；鞍点：10,15 ；极大值 i=3（消亡）：13,22,24
\node[mapdot, draw=blue!55!black]  at (3.6,0)  {0};
\node[mapdot, draw=blue!55!black]  at (3.6,1)  {1};
\node[mapdot, draw=blue!55!black]  at (3.6,8)  {8};
\node[mapdot, draw=green!45!black] at (3.6,10) {10};
\node[mapdot, draw=red!70!black]   at (3.6,13) {13};
\node[mapdot, draw=green!45!black] at (3.6,15) {15};
\node[mapdot, draw=red!70!black]   at (3.6,22) {22};
\node[mapdot, draw=red!70!black]   at (3.6,24) {24};


% 中部标题
\node[font=\sffamily\bfseries\small, text=black!70, align=center] at (4.2,26.9)
   {映射 (map)\\{\footnotesize\mdseries 临界点 $\to$ 标量轴分区}};

% =========================================================
% ===== 中 -> 右：映射箭头（指向已移除的等值面，已一并去除）=====
% =========================================================
\iffalse
\draw[-{Stealth[length=2.4mm, width=2mm]}, line width=0.9pt, draw=blue!55!black]
      (4.85,9)  .. controls (5.3,9)  and (5.3,5)  .. (5.65,5);
\draw[-{Stealth[length=2.4mm, width=2mm]}, line width=0.9pt, draw=green!45!black]
      (4.85,10) .. controls (5.3,10) and (5.3,11) .. (5.65,11);
\draw[-{Stealth[length=2.4mm, width=2mm]}, line width=0.9pt, draw=blue!55!black]
      (4.85,14) .. controls (5.3,14) and (5.3,17) .. (5.65,17);
\draw[-{Stealth[length=2.4mm, width=2mm]}, line width=0.9pt, draw=green!45!black]
      (4.85,15) .. controls (5.3,15) and (5.3,23) .. (5.65,23);
\draw[-{Stealth[length=2.4mm, width=2mm]}, line width=0.9pt, draw=blue!55!black]
      (4.85,23) .. controls (5.3,23) and (5.3,29) .. (5.65,29);
% 层级 / 鞍点标签（箭头起点处）
\node[tagt] at (5.02,9.0)  {$t_1$};
\node[tags] at (5.02,10.0) {$S_2$};
\node[tagt] at (5.02,14.0) {$t_2$};
\node[tags] at (5.02,15.0) {$S_1$};
\node[tagt] at (5.02,23.0) {$t_3$};
\fi

% =========================================================
% ===== 右部：3D 体数据等值面（连通分量 + 亏格；保持不变）=====
% =========================================================
% ===== 右部 3D 等值面内容已移除，仅保留占位区域（位置不变）=====
% （原 blobA/B/C、mergeiso 及 t1/t2/t3、S1/S2 层级已删除；
%   保留右侧画布宽度与标量方向箭头以维持整体布局。）
\iffalse
% blobA：横向拉长、右侧凸起
\newcommand{\blobA}[3]{%
  \begin{scope}[shift={(#1,#2)}, scale=#3, x=1cm, y=1cm]
    \draw[iso3d] plot[smooth cycle, tension=0.75]
      coordinates {(0,0.55)(0.5,0.5)(0.85,0.1)(0.6,-0.28)
                   (0.7,-0.5)(0.1,-0.55)(-0.4,-0.42)(-0.72,-0.05)
                   (-0.5,0.35)};
    \draw[iso3d] plot[smooth cycle, tension=0.8]
      coordinates {(0.12,0.14)(0.28,0.0)(0.16,-0.16)(-0.06,-0.12)(-0.1,0.06)};
  \end{scope}
}
% blobB：竖向、左上凹陷
\newcommand{\blobB}[3]{%
  \begin{scope}[shift={(#1,#2)}, scale=#3, x=1cm, y=1cm]
    \draw[iso3d] plot[smooth cycle, tension=0.7]
      coordinates {(0,0.78)(0.35,0.5)(0.5,0.05)(0.4,-0.45)
                   (0.05,-0.7)(-0.4,-0.5)(-0.55,-0.05)(-0.28,0.3)(-0.5,0.62)};
    \draw[iso3d] plot[smooth cycle, tension=0.8]
      coordinates {(-0.02,0.22)(0.16,0.1)(0.12,-0.14)(-0.12,-0.16)(-0.18,0.04)};
  \end{scope}
}
% blobC：多凸角、不规则
\newcommand{\blobC}[3]{%
  \begin{scope}[shift={(#1,#2)}, scale=#3, x=1cm, y=1cm]
    \draw[iso3d] plot[smooth cycle, tension=0.72]
      coordinates {(0,0.68)(0.55,0.42)(0.42,0.0)(0.72,-0.35)
                   (0.2,-0.62)(-0.28,-0.55)(-0.62,-0.2)(-0.45,0.25)(-0.6,0.55)};
    \draw[iso3d] plot[smooth cycle, tension=0.8]
      coordinates {(0.04,0.18)(0.22,0.06)(0.18,-0.16)(-0.06,-0.2)(-0.2,-0.02)(-0.12,0.14)};
  \end{scope}
}

% ---- t1 层级（槽位 y=5）：c=3 ----
\begin{scope}[shift={(6.7,5)}, x=1cm, y=1cm]
  \node[font=\sffamily\footnotesize\bfseries, text=red!70!black] at (-1.55,0) {$t_1$};
  \blobA{-1.05}{0.15}{0.42}
  \blobB{0.15}{0.5}{0.40}
  \blobC{0.35}{-0.55}{0.45}
  \node[font=\sffamily\footnotesize, text=black!70, align=center] at (-0.2,-1.15)
     {$c=3,\ \chi=6$：三个球面分量};
\end{scope}

% ---- t2 层级（槽位 y=17）：c=2 ----
\begin{scope}[shift={(6.7,17)}, x=1cm, y=1cm]
  \node[font=\sffamily\footnotesize\bfseries, text=red!70!black] at (-1.55,0) {$t_2$};
  \blobB{-0.5}{0.05}{0.72}
  \blobA{0.75}{-0.35}{0.5}
  \node[font=\sffamily\footnotesize, text=black!70, align=center] at (0,-1.15)
     {$c=2,\ \chi=4$：一分量封顶消失};
\end{scope}

% ---- t3 层级（槽位 y=29）：c=1 ----
\begin{scope}[shift={(6.7,29)}, x=1cm, y=1cm]
  \node[font=\sffamily\footnotesize\bfseries, text=red!70!black] at (-1.55,0) {$t_3$};
  \blobC{0}{0}{1.05}
  \node[font=\sffamily\footnotesize, text=black!70, align=center] at (0,-1.15)
     {$c=1,\ \chi=2$：仅剩单一分量};
\end{scope}

% ---- 鞍点合并瞬间的等值面 ----
\newcommand{\mergeiso}[3]{% x y scale
  \begin{scope}[shift={(#1,#2)}, scale=#3, x=1cm, y=1cm]
    \draw[iso3d] plot[smooth cycle, tension=0.7]
      coordinates {(0,0.16)(-0.32,0.42)(-0.78,0.35)(-0.92,-0.05)
                   (-0.72,-0.42)(-0.28,-0.4)(0,-0.16)(-0.06,0)};
    \draw[iso3d] plot[smooth cycle, tension=0.7]
      coordinates {(0,0.16)(0.3,0.4)(0.7,0.36)(0.9,0.0)
                   (0.72,-0.44)(0.26,-0.42)(0,-0.16)(0.06,0)};
    \node[circle, fill=red!80!black, draw=white, line width=0.4pt,
          inner sep=0pt, minimum size=3.4pt] at (0,0) {};
    \draw[iso3d] (-0.48,-0.02) ellipse (0.14 and 0.10);
    \draw[iso3d] ( 0.48,-0.02) ellipse (0.14 and 0.10);
  \end{scope}
}

% ---- S1 = 15 层级：两分量在鞍点合并 ----
\begin{scope}[shift={(6.7,23)}, x=1cm, y=1cm]
  \node[font=\sffamily\footnotesize\bfseries, text=green!45!black] at (-1.55,0) {$S_1$};
  \mergeiso{0}{0}{0.9}
  \node[font=\sffamily\footnotesize, text=black!70, align=center] at (0,-1.1)
     {鞍点 $S_1$（$f\!=\!15$）：两分量相切合并};
\end{scope}

% ---- S2 = 10 层级：两分量在鞍点合并，另有一个独立分量 ----
\begin{scope}[shift={(6.7,11)}, x=1cm, y=1cm]
  \node[font=\sffamily\footnotesize\bfseries, text=green!45!black] at (-1.55,0) {$S_2$};
  \mergeiso{-0.35}{0.15}{0.72}
  \blobA{0.95}{-0.15}{0.42}
  \node[font=\sffamily\footnotesize, text=black!70, align=center] at (0,-1.05)
     {鞍点 $S_2$（$f\!=\!10$）：两分量合并 $+$ 一独立分量};
\end{scope}

% ===== 标量方向箭头 =====
\fi
\draw[-{Stealth[length=2.6mm]}, line width=1pt, black!55] (8.3,0) -- (8.3,24);
\node[font=\sffamily\small, text=black!60, rotate=-90] at (8.55,12) {标量 $f$ 增大};
% 顶部占位（维持原画布上边界，避免整体版面收缩）
\path (6.7,31.5);
% ---- 右部标题 ----
\node[font=\sffamily\bfseries\small, text=black!70, align=center] at (6.9,26.9)
   {真实网格中等值面};

% =========================================================
% ===== 从鞍点 f=10 / f=15 引导线到右侧占位框（等值面图待粘贴）=====
% 起点取色带右边界，终点连到占位框左边缘。
% =========================================================
% ---- 鞍点 S1（f=15）-> 上方占位框 ----
\draw[-{Stealth[length=2.4mm, width=2mm]}, line width=0.9pt, draw=green!45!black]
      (4.85,15) .. controls (5.4,15) and (5.4,20) .. (5.7,20);
\node[draw=green!45!black, dashed, line width=0.8pt, fill=green!4,
      minimum width=3.2cm, minimum height=3.2cm, anchor=south west,
      font=\sffamily\scriptsize, text=black!45, align=center] at (5.7,17)
   {鞍点 $S_1$（$f{=}15$）\\[2pt]等值面待粘贴};

% ---- 鞍点 S2（f=10）-> 下方占位框 ----
\draw[-{Stealth[length=2.4mm, width=2mm]}, line width=0.9pt, draw=green!45!black]
      (4.85,10) .. controls (5.4,10) and (5.4,8) .. (5.7,8);
\node[draw=green!45!black, dashed, line width=0.8pt, fill=green!4,
      minimum width=3.2cm, minimum height=3.2cm, anchor=north west,
      font=\sffamily\scriptsize, text=black!45, align=center] at (5.7,11)
   {鞍点 $S_2$（$f{=}10$）\\[2pt]等值面待粘贴};

% =========================================================
% ===== 底部图例、表名与表注（统一放入一个居中节点，避免坐标缩放导致堆叠）=====
% =========================================================
\node[anchor=north, align=center] at (3.7,-1.2) {%
  \begin{minipage}{15cm}
    \centering
    % ---- 图例 ----
    {\sffamily\footnotesize
      \textcolor{blue!55!black}{$\bullet$}~极小值：分量出生\hspace{2.2em}
      \textcolor{green!45!black}{$\bullet$}~鞍点：合并/分裂\hspace{2.2em}
      \textcolor{red!70!black}{$\bullet$}~极大值：分量消亡\par}
    \vspace{6pt}
    % ---- 区间色带颜色图例（按连通分量数 c）----
    {\sffamily\footnotesize
      区间颜色（按连通分量数 $c$）：\hspace{0.6em}
      \tikz[baseline=-0.5ex]\fill[blue!12]   (0,0) rectangle (0.42,0.26);~$c{=}1$\hspace{1.8em}
      \tikz[baseline=-0.5ex]\fill[green!14]  (0,0) rectangle (0.42,0.26);~$c{=}2$\hspace{1.8em}
      \tikz[baseline=-0.5ex]\fill[orange!20] (0,0) rectangle (0.42,0.26);~$c{=}3$\par}
    \vspace{10pt}
    % ---- 表名 ----
    {\sffamily\bfseries\normalsize 图 1\quad 增广轮廓树 $\to$ 临界值分区映射\par}
    \vspace{3pt}
    {\sffamily\footnotesize\color{black!60}（左：轮廓树；\ 中：临界点到标量轴的映射与区间标注；\ 右：在鞍点 $f{=}10,15$ 处提取的真实网格等值面，展示分量合并瞬间）\par}
    \vspace{10pt}
    % ---- 符号含义说明（表注）----
    {\sffamily\scriptsize\color{black!80}
      \renewcommand{\arraystretch}{1.3}%
      \begin{tabular}{@{}r@{~}l@{\hspace{2.4em}}r@{~}l@{}}
        $f$        & 标量场（等值函数），沿箭头方向递增   &
        $c$        & 区间内等值面的连通分量数 \\
        $\chi$     & 边 / 区间等值面的 Euler 示性数     &
        $I_k$      & 临界值划分出的第 $k$ 个标量区间 \\
        $g$        & 亏格，$g=c-x/2$（$x$ 为每条边 $\chi$，取常数） & & \\
      \end{tabular}\par}
    \vspace{3pt}
    {\sffamily\scriptsize\color{black!60}
      设初始区间 $I_1$ 亏格为 $g_1$，则各区间 $g=g_1+(c-1)$。\par}
  \end{minipage}%
};

\end{tikzpicture}
\end{document}
